\chapter{Text}\label{ch-text}

Text in Zeno is a \verb|str|: a row of characters. It is written
between double quotes, or between single quotes when it holds double
ones:

\begin{zeno}
let name = "Ada"
let quote = 'She said "hello"'
print(name, quote, "it's fine")
\end{zeno}

A string must close on the line it opens on.

\section{Characters that cannot be typed}

A backslash starts an \emph{escape}: two characters that stand for one.

\begin{center}
\begin{tabular}{ll}
\toprule
Escape & Means \\
\midrule
\verb|\n| & a line end \\
\verb|\t| & a tab \\
\verb|\r| & a carriage return \\
\verb|\"|, \verb|\'| & a quote \\
\verb|\\| & a backslash \\
\bottomrule
\end{tabular}
\end{center}

\begin{zeno}
print("first line\nsecond line")
print("a backslash: \\, and a quote: \"")
\end{zeno}

Any other letter after a backslash is a mistake, so a slip such as
\verb|"C:\temp\new"| is caught rather than quietly turned into a tab
and a line end.

\section{Characters, not bytes}

Zeno text is Unicode, kept in UTF-8. Its length, its positions and a
loop over it all count \emph{characters}, whatever alphabet they are
from and however many bytes each one takes:

\begin{zeno}
let word = "naïve"
print(len(word), word[2], word.find("v"))
\end{zeno}

A position counts from 0, so \verb|word[0]| is the first character. A
position past the end, or below 0, is an \verb|IndexError|:

\begin{zeno}[fails=yes]
print("abc"[3])
\end{zeno}

\begin{why}
In Python, \verb|text[-1]| is the last character. That is handy, but
it means a position that goes below zero \emph{by mistake} quietly
reads from the other end instead of failing. A clear error is the
safer thing, above all in code that runs on hardware. (Decision 2.6.)
\end{why}

Text cannot be changed inside. Every method that seems to change it
gives back new text and leaves the old as it was:

\begin{zeno}[fails=yes]
let word = "cat"
word[0] = "b"
\end{zeno}

\section{Joining and comparing}

\verb|+| joins two pieces of text. Both sides must be text; a number is
made into text with \verb|str|, or put into a formatted string:

\begin{zeno}
let count = 3
print("count: " + str(count))
print("ab" == "ab", "ab" < "b", "B" < "a")
print("ell" in "hello", "x" not in "hello")
\end{zeno}

Text is ordered character by character, by the number Unicode gives
each character, so every capital comes before every small letter.
\verb|in| asks whether one piece of text is found in another.

\section{Formatted strings}

An \verb|f| in front of a string lets it hold expressions in braces.
Each one is worked out and its value put into the text, as
\verb|print| would show it:

\begin{zeno}
let name = "Ada"
let items = 3
print(f"{name} has {items} items, {items * 2} halves")
print(f"Shout it: {name.upper()}")
print(f"Braces themselves: {{ and }}")
\end{zeno}

Inside the braces, a string must use the other kind of quote from the
one around it.

\begin{why}
Without the \verb|f|, any string with a \verb|{| in it would change
meaning, and such strings are common: JSON, templates, messages for a
device. The one letter marks the strings that hold code.
(Decision 1.14.)
\end{why}

\subsection{Formats}

After the expression, a \verb|:| and a \emph{format} say how to lay the
value out:

\begin{center}
\begin{tabular}{lll}
\toprule
Format & Means & \verb|f"[{v:...}]"| \\
\midrule
\verb|.2f| & two digits after the point & \verb|[3.14]| for \verb|3.14159| \\
\verb|8.3f| & at least 8 wide, 3 digits after the point & \verb|[   3.142]| \\
\verb|6| & at least 6 wide; numbers on the right, text on the left & \verb|[    42]| \\
\verb|<6|, \verb|>6|, \verb|^6| & at least 6 wide: left, right, centred & \verb|[42    ]| \\
\verb|03| & at least 3 wide, with zeros in front & \verb|[042]| \\
\verb|x|, \verb|X| & a whole number in hexadecimal & \verb|[ff]|, \verb|[FF]| \\
\verb|b| & a whole number in binary & \verb|[101]| \\
\verb|d| & a whole number & \verb|[42]| \\
\bottomrule
\end{tabular}
\end{center}

\begin{zeno}
let pi = 3.14159
let n = 42
print(f"[{pi:.2f}] [{pi:8.3f}] [{n:6}] [{n:<6}] [{n:^6}] [{n:04}]")
print(f"[{255:x}] [{255:X}] [{5:b}] [{0xA5:08b}]")
print(f"[{'ab':>5}] [{'ab':<5}]")
\end{zeno}

A format that makes no sense for the value is refused before the
program runs:

\begin{zeno}[fails=yes]
let pi = 3.14159
print(f"{pi:x}")
\end{zeno}

\section{The methods of text}

A \emph{method} is a function that belongs to a value, and is called
with a dot after it: \verb|name.upper()|. Text has these:

\begin{center}
\begin{tabular}{lp{8.5cm}}
\toprule
Method & Gives \\
\midrule
\verb|upper()|, \verb|lower()| & the text in capitals, or in small letters \\
\verb|strip()| & the text without blanks at either end \\
\verb|split()| & a list of the words, split at blanks \\
\verb|split(sep)| & a list of the pieces between each \verb|sep| \\
\verb|lines()| & a list of the lines, without their line ends \\
\verb|replace(old, new)| & the text with every \verb|old| made \verb|new| \\
\verb|starts_with(part)|, \verb|ends_with(part)| & whether it starts or ends with \verb|part| \\
\verb|find(part)| & the position where \verb|part| starts, or \verb|none| \\
\verb|slice(from, to)| & the characters from \verb|from| up to, not including, \verb|to| \\
\verb|slice(from)| & the characters from \verb|from| to the end \\
\verb|code()| & the Unicode number of a text of one character \\
\bottomrule
\end{tabular}
\end{center}

\begin{zeno}
let line = "  Set Speed 40  "
let words = line.strip().lower().split()
print(words)
print("a,b,,c".split(","), "x=1".replace("=", " is "))
print("hello".starts_with("he"), "hello".ends_with("x"))
print("hello".slice(1, 4), "hello".slice(2))
print("A".code(), char(66), char(8364))
\end{zeno}

\verb|char(n)| is the other way round from \verb|code()|: the
character with that number. It is a function and not a method because
there is no text yet to ask: it \emph{makes} text from a number, as
\verb|str(n)| does.

\subsection{split and lines}

\verb|split()| with nothing in the brackets splits at any run of blanks
and never gives an empty piece. \verb|split(sep)| splits at each
\verb|sep| exactly, so two in a row give an empty piece between them:

\begin{zeno}
print("one  two   three".split())
print("  a b ".split(" "))
print("".split(","), "a,".split(","))
\end{zeno}

\verb|lines()| is the way to cut text into lines. It knows both kinds
of line end, \verb|"\n"| and \verb|"\r\n"|, and a last line with no end:

\begin{zeno}
print("a\nb\n".split("\n"))
print("a\nb\n".lines())
print("one\r\ntwo".lines())
\end{zeno}

\subsection{find gives a position or none}

\verb|find| gives the position of a piece of text, or \verb|none| when
it is not there. Because the answer may be \verb|none|, Zeno will not
let it be used as a number until the program has asked:

\begin{zeno}
function first_word(line: str) -> str {
    let space = line.find(" ")
    if space is none {
        return line
    }
    return line.slice(0, space)
}

print(first_word("set speed 40"), first_word("stop"))
\end{zeno}

\begin{why}
Python's \verb|find| gives \verb|-1| when the text is not there, and
\verb|-1| is a number, so nothing stops it from being used as a
position. \verb|none| cannot be used as a position at all, so the
mistake is caught before the program runs. (Decision 6.3.)
\end{why}

\hyperref[ch-none]{Values That May Be Missing} is about \verb|none| and
how Zeno knows when it has been ruled out.

\subsection{slice takes a part}

\verb|slice(from, to)| counts as \verb|range| does: from \verb|from| up
to but not including \verb|to|. A position outside the text is an
error, not a shorter slice:

\begin{zeno}[fails=yes]
print("abc".slice(1, 9))
\end{zeno}

If you write Python's \verb|text[1:3]|, Zeno says what its own way is:

\begin{zeno}[fails=yes]
let text = "hello"
print(text[1:3])
\end{zeno}

\subsection{Plain letters only}

\verb|upper| and \verb|lower| change the letters \verb|a| to \verb|z|
and \verb|A| to \verb|Z| and nothing else, and \verb|strip| and
\verb|split()| know the space, the tab, the line ends and the form
feed as blanks. Letters of other alphabets pass through untouched:

\begin{zeno}
print("Straße".upper(), "ÉCOLE".lower())
\end{zeno}

\begin{why}
The full rules for every alphabet need tables of tens of kilobytes,
and a compiled Zeno program has to fit on a microcontroller. A method
must also give the same answer under \verb|zeno run| as in a compiled
program, so the rule is the same everywhere. (Decision 3.14.)
\end{why}

\section{A loop over text}

A \verb|for| loop over text is given each character in turn. With two
names it is given each position and the character there:

\begin{zeno}
for ch in "abc" {
    print(ch)
}
for at, ch in "hé" {
    print(at, ch)
}
\end{zeno}

\hyperref[ch-control-flow]{Deciding and Repeating} has loops in full.

\section{Text from other values}

\verb|str(x)| gives any value as text, exactly as \verb|print| shows
it. Inside a list, text is shown in quotes, so that \verb|["1"]| and
\verb|[1]| look different:

\begin{zeno}
print(str(42) + "!", str([1, "a"]), str(none), str(true))
\end{zeno}

\begin{summary}
A \verb|str| is text, counted in characters. It is written in quotes,
with escapes after a backslash, and never changes inside. \verb|+|
joins text, \verb|in| finds it, and an \verb|f| string puts values into
it, laid out by a format. Its methods give new text, a list of pieces,
a position or \verb|none|.
\end{summary}

\begin{tryit}
\begin{enumerate}
\item Write \verb|function last_word(line: str) -> str|. There is no
  method that searches from the end. What is the shortest way with what
  there is?
\item Print a table of the numbers 1 to 5 with their squares and cubes,
  each column 6 wide, using formats.
\item Print \verb|"€".code()| and work out by hand, from the number,
  how many bytes the character takes in UTF-8.
\item What does \verb|"a,b,c".split(",")| give? And
  \verb|",".split(",")|? Say why each piece is there.
\item Write \verb|"C:\temp"| in a program and read the message.
\end{enumerate}
\end{tryit}
